\begin{trapbox}
\begin{itemize}
\item \textbf{把内角平分线写成外角平分线：}
双曲线在 $P$ 处的\emph{切线}平分 $\angle F_1PF_2$ 的\textbf{内角}；
同一对焦点射线的外角平分线与切线垂直，是法线方向。
椭圆的切线/法线对应关系恰好与此不同。

\item \textbf{把右支公式机械照搬到左支：}
左支 $x_0<-a$，对应 $m<0$。安全的统一写法是
$\boxed{x_0m=a^2}$，或在证明中保留 $|r_1-r_2|=2a$，最后平方。

\item \textbf{混淆双曲线与椭圆焦点关系：}
本题 $c^2=a^2+b^2$，并非 $c^2=a^2-b^2$；且 $c>a$。

\item \textbf{角平分线的比值方向写反：}
$\dfrac{F_1M}{MF_2}=\dfrac{PF_1}{PF_2}$，
对应为\emph{左底段/右底段 = 左焦距/右焦距}，
不能无缘无故倒置一侧的次序。

\item \textbf{直接写 $F_1M=c+m$ 却忘记验证 $M$ 的位置：}
真正的\emph{内角}平分线交点在\emph{开线段} $F_1F_2$ 内，故 $-c<m<c$。
本结论还进一步得到 $|m|<a<c$。

\item \textbf{把内切圆与焦点轴的切点 $T$ 认成 $M$：}
$T=(\operatorname{sgn}x_0)a$（分支顶点），
而 $M=(a^2/x_0,0)$ 随 $x_0$ 变化，且非顶点时 $T\ne M$。
内切圆方法能解释固定切点，不意味着\emph{角平分线必经过切点}。

\item \textbf{忘记排除顶点：}
在 $P=(\pm a,0)$ 时 $P,F_1,F_2$ 共线，所谓焦点三角形退化，
从 $P$ 到轴交点 $M$ 无法按题意唯一确定；不能拿退化情形套角平分线定理。
但顶点切线 $x=\pm a$ 仍有定义。

\item \textbf{把“最快引用”误作独立证明：}
若直接说“内角平分线就是切线”，再由切线求截距，
则已假设了结论二；独立证明应先使用方法一，并借双曲线方程反推出切线方程。

\item \textbf{混淆 $P$ 和 $M$ 的横坐标：}
$P=(x_0,y_0)$ 在双曲线上，$M=(m,0)$ 在焦点轴上；
两者满足\emph{乘积} $x_0m=a^2$，不是 $m=x_0$ 或 $m=a^2x_0$。
\end{itemize}
\end{trapbox}
